
\paragraph{An algebraic vector and the Galois comparison.}
Ramsey's section model \cite[Definition~3.1 and Theorem~4.2]{ramsey} realizes the half-integral-weight space as
sections of the indicated divisor sheaf on $X_1(4N_{\rm aux})$.
Here $N_{\rm aux}$ is supported on $S_0$ and includes the chosen
local depths; it need not equal the conductor $N$ of $E$.
Over characteristic-zero fields, flat field extension commutes
with these sections. Cuspidality, scalar character, support
conditions and prescribed good-Hecke eigenvalues are algebraic
linear conditions. Although infinitely many coefficients or
primes can occur in their description, a finite subset cuts out
the same subspace, by finite dimensionality.

One need not impose Petersson orthogonality as an additional
transcendental linear condition. The elliptic good-Hecke eigenspace
is automatically orthogonal to unary theta series: at a good prime,
the unary eigenvalues have absolute value $p+1$, whereas
$|\eta(p)a_p(E)|\le2\sqrt p<p+1$. The unary space and its
orthogonal complement are stable under the Hecke operators and
their adjoints. Thus their elliptic eigenspaces are disjoint.
The nonzero reference functional on the complex space is consequently
nonzero on its algebraic model. Choose such an algebraic $g$ and
scale it to $a(D_0)=1$.
This uses flat field extension only, not Ramsey's arbitrary-ring
base-change assertion in weight $3/2$; see \cite[Remark~4.4]{ramsey}.
Here the field-extension assertion can be checked without a
cohomology-vanishing theorem. Work over a number field $K$ containing
the character values and the roots of unity specifying the cusps and
the support idempotents. A finite affine cover of the modular curve
computes global sections as the kernel of the difference of restriction
maps to pairwise intersections. Tensoring with a field extension is
flat and preserves that kernel. Thus the sections in Ramsey's
Definition~3.1 commute with $K\subset\overline\Q\subset\mathbf C$,
also in weight $3/2$. Impose vanishing at the cusps, the diamond
character, the support idempotents, and the kernels of
$T(p^2)-\eta(p)a_p(E)$ inside this finite-dimensional space.
The span in its dual of all the defining linear equations has a finite
basis. These conditions therefore commute with the same field
extensions. Since the reference functional is nonzero after extension
to $\mathbf C$, it is nonzero before extension; a vector in a $K$-basis has
nonzero reference value and may be divided by that value.
This also places every coefficient of the chosen form in one number
field. No assertion about base change to a residue field, nor about
vanishing of $H^1$ in weight $3/2$, is used.

